\documentclass[11pt,a4paper]{article}
\usepackage[margin=2.2cm]{geometry}
\usepackage[T1]{fontenc}
\usepackage{charter}
\usepackage{microtype}
\usepackage{graphicx}
\usepackage{booktabs}
\usepackage{array}
\usepackage{tabularx}
\usepackage{float}
\usepackage{caption}
\usepackage{enumitem}
\usepackage{amsmath}
\usepackage[table]{xcolor}
\usepackage{listings}
\usepackage[most]{tcolorbox}
\usepackage{fancyhdr}
\usepackage[colorlinks=true,linkcolor=accent,urlcolor=accent]{hyperref}

\definecolor{accent}{HTML}{3C3489}
\definecolor{codebg}{HTML}{F6F6F3}
\definecolor{kw}{HTML}{185FA5}
\definecolor{cm}{HTML}{6B6A64}
\definecolor{okgreen}{HTML}{0F6E56}
\definecolor{badred}{HTML}{A32D2D}

\setlength{\parindent}{0pt}
\setlength{\parskip}{6pt}
\setlist{itemsep=2pt,topsep=4pt}
\captionsetup{font=small,labelfont=bf}
\renewcommand{\arraystretch}{1.18}

\pagestyle{fancy}\fancyhf{}
\rhead{\small\color{cm}\reporttag}
\cfoot{\small\thepage}
\renewcommand{\headrulewidth}{0.3pt}

\lstdefinelanguage{SV}{
  morekeywords={module,endmodule,input,output,inout,logic,wire,reg,always_ff,always_comb,assign,
    begin,end,if,else,case,endcase,default,typedef,enum,localparam,posedge,negedge,or,task,endtask,
    function,endfunction,automatic,initial,repeat,while,for,integer,string,signed,void},
  sensitive=true, morecomment=[l]{//}, morecomment=[s]{/*}{*/}, morestring=[b]"}
\lstdefinelanguage{DCTcl}{
  morekeywords={set,lappend,define_design_lib,analyze,elaborate,current_design,link,check_design,
    create_clock,set_input_delay,set_output_delay,set_max_area,compile,report_area,report_timing,
    report_qor,report_resources,report_constraints,write,write_sdc,set_dont_use,get_lib_cells,puts,gui_start},
  sensitive=true, morecomment=[l]{\#}}
\lstset{basicstyle=\ttfamily\footnotesize, backgroundcolor=\color{codebg}, frame=single,
  rulecolor=\color{codebg}, keywordstyle=\color{kw}\bfseries, commentstyle=\color{cm}\itshape,
  columns=fullflexible, keepspaces=true, breaklines=true, showstringspaces=false,
  xleftmargin=4pt, xrightmargin=4pt, aboveskip=6pt, belowskip=6pt}

\newtcolorbox{task}{colback=accent!5,colframe=accent,boxrule=0.6pt,arc=2pt,
  left=8pt,right=8pt,top=4pt,bottom=4pt,fonttitle=\bfseries,title=The assignment}
\newtcolorbox{keypoint}[1]{colback=okgreen!6,colframe=okgreen,boxrule=0.6pt,arc=2pt,
  left=8pt,right=8pt,top=4pt,bottom=4pt,fonttitle=\bfseries,title={#1}}
\newtcolorbox{lesson}[1]{colback=orange!7,colframe=orange!70!black,boxrule=0.6pt,arc=2pt,
  left=8pt,right=8pt,top=4pt,bottom=4pt,fonttitle=\bfseries,title={#1}}

\newcommand{\fig}[3]{\begin{figure}[H]\centering\includegraphics[width=#2\linewidth]{figures/#1}\caption{#3}\end{figure}}
\newcommand{\pass}{\textcolor{okgreen}{\textbf{PASS}}}
\newcommand{\met}{\textcolor{okgreen}{\textbf{MET}}}
\newcommand{\viol}{\textcolor{badred}{\textbf{VIOLATED}}}

\newcommand{\header}[3]{%
  \begin{center}
  {\small\color{cm} CPE 527 VLSI Design \quad|\quad Lab 2 \quad|\quad University of Alabama in Huntsville}\\[8pt]
  {\LARGE\bfseries\color{accent} #1}\\[4pt]
  {\large #2}\\[8pt]
  {\small Bibek Shrestha \quad|\quad September 2026 \quad|\quad #3}
  \end{center}\vspace{2pt}\hrule\vspace{6pt}}
\newcommand{\reporttag}{Lab 2 / Assignment 3: MSP430 hardware multiplier}
\begin{document}
\header{MSP430 Hardware Multiplier Accelerator}{MPY, MPYS, MAC and MACS over a 16-bit bus}{SystemVerilog, VCS, Design Compiler, SAED 90\,nm}

\begin{task}
Design and test a hardware multiplier for the MSP430 that supports unsigned (MPY) and signed (MPYS) multiply,
and unsigned (MAC) and signed (MACS) multiply-accumulate. The CPU picks the mode and first operand by writing to
0x0130 / 0x0132 / 0x0134 / 0x0136. Writing the second operand to 0x0138 starts the operation. The 32-bit result
is read at 0x013A (low 16 bits) and 0x013C (high 16 bits). Ports: \texttt{clk}, \texttt{reset\_n} (active low),
16-bit \texttt{addr}, 16-bit bidirectional \texttt{data\_bus}, \texttt{write\_enable}, \texttt{read\_enable},
and output \texttt{busy}. Write a testbench that verifies the design thoroughly.
\end{task}

\section*{At a glance}
\begin{table}[H]\centering\small
\begin{tabular}{ll}
\toprule
Method & 16-cycle shift-and-add on magnitudes, sign fixed at the end \\
Flip-flops / cells & 153 / 881 \\
Area (cell / total) & 12677.53 / 13462.00 \\
Clock & 7 ns ($\approx$143 MHz), slack 0.00 ns \met \\
Critical path & \texttt{mcand\_reg[0]} $\rightarrow$ 32-bit adder $\rightarrow$ \texttt{product\_reg[31]} \\
Busy time & 18 clocks per operation \\
Tests & 228 (RTL) and 228 (gate level), \textbf{0 errors} \\
Tri-state check & 0 BSLEX cells, 16 TNBUFFX cells \\
\bottomrule
\end{tabular}
\end{table}

\section{The four modes}
All four multiply two 16-bit numbers into a 32-bit result. They differ in two ways: are the numbers signed,
and does the new product replace the result or add to it?

\begin{table}[H]\centering\small
\begin{tabular}{llllc}
\toprule
Mode & op1 address & Numbers & Result & \texttt{mode[1:0]} \\
\midrule
MPY & 0x0130 & unsigned & result = a $\times$ b & 00 \\
MPYS & 0x0132 & signed & result = a $\times$ b & 01 \\
MAC & 0x0134 & unsigned & result = result + a $\times$ b & 10 \\
MACS & 0x0136 & signed & result = result + a $\times$ b & 11 \\
\midrule
OP2 & 0x0138 & \multicolumn{3}{l}{write op2 and \textbf{start}} \\
RESL / RESH & 0x013A / 0x013C & \multicolumn{3}{l}{read result[15:0] / result[31:16]} \\
\bottomrule
\end{tabular}
\caption{Register map. Mode bit 0 = signed, bit 1 = accumulate (it equals address bits [2:1]).}
\end{table}

Example of why signed matters: \texttt{0xFFFF} is 65535 unsigned but $-1$ signed.
MPY gives $65535 \times 65535 = \texttt{0xFFFE0001}$; MPYS gives $(-1)\times(-1) = \texttt{0x00000001}$.

\section{Design}
\fig{arch.png}{0.85}{Architecture: a bus interface and registers on the CPU side, a control FSM and a datapath on the compute side.}

\begin{table}[H]\centering\small
\begin{tabularx}{\linewidth}{lXl}
\toprule
Module & Job & Registers \\
\midrule
\texttt{mult430} (top) & Address decoder, op1/op2/mode registers, read mux, tri-state \texttt{data\_bus} & 34 \\
\texttt{mult430\_ctrl} & FSM and step counter, drives \texttt{busy} & 6 \\
\texttt{mult430\_dp} & Magnitudes, shift-and-add, sign fix, accumulate & 113 \\
\bottomrule
\end{tabularx}
\caption{Three modules, 153 flip-flops in total.}
\end{table}

\subsection*{How one multiplier handles all four modes}
\begin{enumerate}
\item \textbf{Signed modes:} make negative operands positive ($|x| = \sim x + 1$) and remember if the signs differ.
\item \textbf{Multiply} the two positive numbers with shift-and-add, one bit per clock, 16 clocks.
\item \textbf{Fix the sign:} if the signs differed, negate the product.
\item \textbf{Accumulate modes:} add the product to the old result instead of replacing it.
\end{enumerate}
\fig{a3_algo.png}{0.72}{One operation from start to finish.}
\fig{a3_fsm.png}{0.85}{Control FSM. \texttt{busy} is high in every state except IDLE: 1 + 16 + 1 = 18 clocks.}

\subsection*{Worked example: MPY 3 $\times$ 5}
\begin{table}[H]\centering\small
\begin{tabular}{ccccc}
\toprule
Step & multiplier bit & Action & product & multiplicand after \\
\midrule
1 & 1 & add 3 & 3 & 6 \\
2 & 0 & skip & 3 & 12 \\
3 & 1 & add 12 & 15 & 24 \\
4 to 16 & 0 & skip & 15 & \dots \\
\bottomrule
\end{tabular}
\caption{Shift-and-add. For MPYS $-3 \times 5$ the same steps give 15, then negation gives \texttt{0xFFFFFFF1} = $-15$.}
\end{table}

\subsection*{The shared bus}
The data bus is used by both sides, so only one may drive it at a time. The multiplier drives it only during a
read of 0x013A or 0x013C; otherwise it lets go (\texttt{z}). Writes are ignored while \texttt{busy} is high.
\begin{lstlisting}[language=SV]
assign wr_ok    = write_enable && !read_enable && !busy;
assign start    = wr_ok && (addr == ADDR_OP2);
assign data_bus = rd_en ? rd_data : 16'bz;   // tri-state driver
\end{lstlisting}

\section{Testbench}
The testbench acts as the CPU. \texttt{cpu\_write} and \texttt{cpu\_read} perform bus cycles,
a small software model computes the expected answer, and every result is checked automatically.
Inputs change on the falling clock edge to avoid races.

\begin{table}[H]\centering\small
\begin{tabularx}{\linewidth}{lcXc}
\toprule
Group & Tests & What it checks & \\
\midrule
MPY / MPYS / MAC / MACS & 11 & Each mode, largest values, $-32768 \times -32768$, accumulation up and down & \pass \\
Write while busy & 1 & A MAC write during an MPY is ignored (0x31, not 0x72) & \pass \\
Bus released & 2 & \texttt{data\_bus} is \texttt{z} when idle and on a non-result read & \pass \\
Mode switching & 5 & The mode never ``sticks'' from the previous operation & \pass \\
Overflow & 5 & The 32-bit accumulator wraps around both ways & \pass \\
Reset mid-operation & 4 & \texttt{busy} drops instantly, result cleared, then works again & \pass \\
Random & 200 & All modes, seed 430, biased to 0000/FFFF/8000/7FFF & \pass \\
\midrule
\textbf{Total} & \textbf{228} & & \textbf{0 errors} \\
\bottomrule
\end{tabularx}
\caption{Test plan. The random split was MPY 51, MPYS 48, MAC 53, MACS 48.}
\end{table}

\begin{table}[H]\centering\small
\begin{tabular}{lcc}
\toprule
Operation & Product & Result \\
\midrule
MPYS 0 $\times$ $-1$ & 0 & 0x00000000 \\
MAC 2 $\times$ 10 & +20 & 0x00000014 \\
MAC 4 $\times$ 10 & +40 & 0x0000003C \\
MACS $-2$ $\times$ 10 & $-20$ & 0x00000028 \\
MACS $-5$ $\times$ $-5$ & +25 & 0x00000041 \\
\bottomrule
\end{tabular}
\caption{MAC and MACS keep a running total in the same result register.}
\end{table}

\fig{wave12_mpys.png}{1.0}{MPYS $-3\times5$: \texttt{busy} goes high after the op2 write; the two reads return \texttt{fff1} and \texttt{ffff} = \texttt{0xFFFFFFF1}.}
\fig{wave_full.png}{1.0}{The full run: \texttt{tests} reaches 228 and \texttt{errors} stays 0.}

\section{Synthesis}
\subsection*{Tri-state fix}
The TA warned that Design Compiler may build the tri-state bus from \texttt{BSLEX} cells, which ignore the enable
and act like inverters (3 + 5 would read as 247). Blocking them leaves the correct \texttt{TNBUFFX} cells:
\begin{lstlisting}[language=DCTcl]
set_dont_use [get_lib_cells */BSLEX1]
set_dont_use [get_lib_cells */BSLEX2]
set_dont_use [get_lib_cells */BSLEX4]
\end{lstlisting}

\subsection*{Getting timing to pass}
\begin{table}[H]\centering\small
\begin{tabular}{lccc}
\toprule
 & 5 ns & 5 ns, speed effort & \textbf{7 ns (final)} \\
\midrule
Slack & $-1.29$ \viol & $-1.37$ \viol & \textbf{0.00} \met \\
Violating paths & 19 & 17 & 0 \\
Critical path & 6.21 ns & 6.29 ns & 6.91 ns \\
Cells & 998 & 1019 & 881 \\
Design area & 14542.17 & 14666.44 & 13462.00 \\
\bottomrule
\end{tabular}
\caption{The 32-bit ripple-carry adder needs about 6.3\,ns, so 5\,ns could not pass. The spec gives no clock speed, so I used 7\,ns.}
\end{table}

\fig{a3_budget.png}{0.85}{Critical-path timing at 5\,ns and 7\,ns.}

At 7\,ns the path is slightly \emph{longer} (6.91\,ns) but still passes. With \texttt{set\_max\_area 0}, DC uses
spare time to pick smaller, slower gates, which also cut the area by about 1080 units.

\begin{table}[H]\centering\small
\begin{tabular}{lr}
\toprule
Combinational area & 7521.18 \\
Flip-flop area & 5156.35 \\
Net area & 784.47 \\
\textbf{Total} & \textbf{13462.00} \\
\bottomrule
\end{tabular}
\caption{Area at 7\,ns.}
\end{table}

\begin{keypoint}{Critical path}
\texttt{u\_dp/mcand\_reg[0]} $\rightarrow$ 32-bit ripple-carry adder (\texttt{product + mcand}) $\rightarrow$
\texttt{u\_dp/product\_reg[31]}. 53 logic levels, 6.91\,ns, slack 0.00\,ns at 7\,ns ($\approx$143\,MHz).
\end{keypoint}

\section{Gate-level check}
\begin{lstlisting}[language=bash]
grep -c BSLEX  output/assignment3_synth.v     # 0
grep -c TNBUFF output/assignment3_synth.v     # 16 (one per data_bus bit)
./ps_sim/assignment3_synth.simv               # 228 tests, 0 errors
\end{lstlisting}
The netlist passes the same 228 tests with the same random split. The bus-release tests passing on real gates
shows the tri-state fix worked.

\section{Design Vision}
\fig{a3_dv_symbol.png}{0.6}{The synthesized block has exactly the ports in the spec; \texttt{data\_bus} shows the two-way arrow of an \texttt{inout}.}
\fig{a3_dv_path.png}{0.95}{The critical path: the staircase is the 32-bit carry chain.}
\fig{dv_paths.png}{0.95}{Slack of the 50 worst paths: none negative. The worst ones end at \texttt{product\_reg[31]} and \texttt{result\_reg[31]}.}

\section{Choices I made}
\begin{table}[H]\centering\small
\begin{tabularx}{\linewidth}{lX}
\toprule
Clock & 7\,ns, because 5\,ns failed and the spec sets no speed \\
Writes during busy & Ignored, so an operation cannot be corrupted \\
Accumulator & One shared 32-bit result; MPY/MPYS overwrite, MAC/MACS add \\
Overflow & Wraps at 32 bits (no overflow flag in the spec) \\
Data bus & Real \texttt{inout} with the BSLEX fix \\
\bottomrule
\end{tabularx}
\end{table}

\section{What I learned}
\begin{itemize}
\item One unsigned multiplier plus a sign fix covers all four modes.
\item A wide adder sets the clock speed: 32-bit ripple carry limited the design to about 143\,MHz.
      Adding 16 bits per step instead would be the next improvement.
\item A test is only useful if it can fail: the write-while-busy test writes to the MAC address for that reason.
\item Always check the netlist for known bad cells (BSLEX) and re-run the full testbench on the gates.
\end{itemize}

\appendix
\section{Source files}
\subsection*{RTL: \texttt{src/assignment3.sv}}
\lstinputlisting[language=SV]{src/assignment3.sv}
\subsection*{Testbench: \texttt{src/assignment3\_tb.sv}}
\lstinputlisting[language=SV]{src/assignment3_tb.sv}
\subsection*{Synthesis script: \texttt{src/syn.tcl}}
\lstinputlisting[language=DCTcl]{src/syn.tcl}
\end{document}
